\documentclass[10pt,a4paper]{amsart}
\usepackage[margin=19mm]{geometry}
\usepackage{amsmath,amssymb,mathtools}
\usepackage[scr=rsfs,cal=euler]{mathalfa}
\usepackage{xfrac}
\usepackage[unicode,hidelinks,bookmarks=false]{hyperref}
\newcommand{\ie}{i.e.\ }
\newcommand{\cf}{cf.\ }
\newcommand{\resp}{respectively\ }
\newcommand{\Subsection}{\subsection}
\providecommand{\nobreakdash}{\nobreak\hbox{-}}
\setlength{\emergencystretch}{2em}
\multlinegap0pt
\usepackage{bm}
\usepackage{bbm}
\usepackage{dsfont}
\usepackage{xspace}
\usepackage{cancel}
\usepackage{mathtools}
\usepackage{amsthm}
\makeatletter
\@ifundefined{thm}{\newtheorem{thm}{Thm}}{}
\@ifundefined{lem}{\newtheorem{lem}[thm]{Lemma}}{}
\makeatother
\newcommand{\C}{\mathbb{C}}
\newcommand{\M}{\mathrm{M}}
\newcommand{\cA}{\mathcal{A}}
\newcommand{\cD}{\mathcal{D}}
\newcommand{\cH}{\mathcal{H}}
\newcommand{\cK}{\mathcal{K}}
\newcommand{\cO}{\mathcal{O}}
\newcommand{\cR}{\mathcal{R}}
\renewcommand\Re{\mathrm{Re}}
\title[Matrix systems, algebras, and open maps]{Matrix systems, algebras, and open maps}
\author{Stephan Weis}
\date{Source: arXiv:2502.01761v3 (CC BY 4.0)}
\begin{document}
\maketitle

\noindent\emph{Source and licence.} Matrix systems, algebras, and open maps. Version arXiv:2502.01761v3 (CC BY 4.0), distributed under the Creative Commons Attribution 4.0 International licence, as recorded in the arXiv licence metadata for this exact version and shown on the abstract page https://arxiv.org/abs/2502.01761v3. Full rights notice: licenses/arxiv-2502.01761-CC-BY-4.0.txt.\par\smallskip

\noindent\textit{Editorial scope.} Counted unit: the Lem whose source label is \texttt{lem:cyclic-group} in the article (source0.tex lines 1409--1415, source label \texttt{lem:cyclic-group}), delivered together with the article's own complete original proof of it (source0.tex lines 1417--1478, one \texttt{proof} environment of 62 lines). No mathematics is written, repaired or re-derived: statement and proof are the article's own text line for line. Required context retained uncounted: eq:D-inter (lines 673--675). Every definition, display and label of those lines is retained unchanged. Omitted from this package and not redistributed: 1--672, 676--1408, 1416--1416, 1479--2795. Nothing inside a retained excerpt is omitted or altered: every retained body is delivered exactly as the article's lines are written, so no omission receipt of any kind accompanies this package. Citations: the retained excerpts carry no citation command at all, so no bibliography entry of the article is transcribed and no reference is redistributed. Reference adaptation: none was needed and none was made; every reference inside the delivered unit resolves inside it, so no unresolved reference or citation is produced. Printed slips kept literal: no printed word, symbol, hypothesis or number of the counted statement or of its proof was altered. Layout and class: the article's own class is \texttt{amsart} with packages including inputenc, fontenc, babel, amsfonts, amsmath, amssymb, dsfont, braket, xy, color, graphicx, hyperref; the wrapper is the lane's pinned \texttt{amsart} editorial wrapper at 10pt on A4, which restates the theorem-like environments the retained text needs and adds the article's own macro definitions that the retained text needs (\texttt{\textbackslash C}, \texttt{\textbackslash M}, \texttt{\textbackslash Re}, \texttt{\textbackslash cA}, \texttt{\textbackslash cD}, \texttt{\textbackslash cH}, \texttt{\textbackslash cK}, \texttt{\textbackslash cO}, \texttt{\textbackslash cR}), with extra wrapper packages (bm, bbm, dsfont, xspace, cancel, mathtools). The delivered numbering is the unit's own. Assets: no retained span contains a figure environment, a graphics-inclusion command, a PDF-page inclusion, a table or a TikZ picture; no image file is copied into this package, the package declares no asset dependency and no image file is redistributed with it. Licence: version arXiv:2502.01761v3 of this article is distributed under the Creative Commons Attribution 4.0 International licence, as recorded in the arXiv licence metadata for this exact version and shown on the abstract page https://arxiv.org/abs/2502.01761v3. A full rights notice is delivered at licenses/arxiv-2502.01761-CC-BY-4.0.txt. Characters: every retained excerpt is pure ASCII, so the delivered engine, XeLaTeX with its default Latin Modern text font, reports no missing character.
\begin{equation}\label{eq:D-inter}
\cD_2=\cD_1\cap\cR_2\,.
\end{equation}

\begin{lem}\label{lem:cyclic-group}
Let $\cR$ be a real matrix system on $\C^n$ and let $\cD(\cR)$ be stable
and invariant under an orthogonal transformation $\gamma:\cR\to\cR$ that 
generates a finite cyclic group. Let $\cA:=\{A\in\cR:\gamma(A)=A\}$ be a 
real *-subalgebra of $\M_n$. Then the orthogonal projection 
$\cD(\cR)\to\cD(\cA)$ is open. 
\end{lem}

\begin{proof}
Let $k$ denote the order of the cyclic group generated by $\gamma$. The main idea 
is to write the orthogonal projection $\pi:\cD(\cR)\to\cD(\cA)$ in terms of the 
arithmetic mean map
\[\textstyle
a:\cD(\cR)^{\times k}\to\cD(\cR)\,,
\quad
(\rho_1,\dots,\rho_k)\mapsto\frac{1}{k}\sum_{i=1}^k\rho_i\,,
\]
which, as discussed above, is open because $\cD(\cR)$ is stable. The proof is done 
in two steps. First, we prove that 
\begin{align}\label{eq:group1}
a(\gamma(\cO)\times\dots\times\gamma^k(\cO))\cap\cA
\end{align}
is an open subset of $\cD(\cA)$ for all open subsets $\cO$ of 
$\cD(\cR)$. Secondly, we prove that 
\begin{align}\label{eq:group2}
a(\gamma(\cK)\times\dots\times\gamma^k(\cK))\cap\cA=\pi(\cK)
\end{align}
holds for all convex subsets $\cK$ of $\cD(\cR)$. The assertions \eqref{eq:group1}
and \eqref{eq:group2} together show that $\pi(\cO)$ is an open subset of $\cD(\cA)$ 
for all convex open subsets $\cO$ of $\cD(\cR)$, and hence for all open subsets.
\par
%
First, we prove that the set in \eqref{eq:group1} is open. As $\cD(\cR)$ is 
invariant under the orthogonal transformation $\gamma$, the map $\gamma$ 
restricts to a homeomorphism $\cD(\cR)\to\cD(\cR)$. It follows that 
$\gamma(\cO)\times\dots\times\gamma^k(\cO)$ is an open subset of the $k$-fold 
cartesian product $\cD(\cR)^{\times k}$. Then
\[
\widetilde{\cO}:=a(\gamma(\cO)\times\dots\times\gamma^k(\cO))
\]
is an open subset of $\cD(\cR)$, because $\cD(\cR)$ is stable. Finally, 
$\widetilde{\cO}\cap\cA$ is an open subset of $\cD(\cA)$ by equation 
\eqref{eq:D-inter}, which shows that $\widetilde{\cO}\cap\cA$ equals 
$\widetilde{\cO}\cap\cD(\cA)$. 
\par
%
Secondly, we prove the formula \eqref{eq:group2}, beginning with the inclusion 
``$\supset$''. Let $\rho\in\cK\subset\cD(\cR)$. The density matrix
$\sigma:=a(\gamma(\rho),\dots,\gamma^k(\rho))$ is invariant under $\gamma$. 
This implies $\sigma\in\cA$, hence $\sigma\in\cD(\cA)$ by \eqref{eq:D-inter}. 
Since $\gamma$ is self-adjoint, for all $A\in\cH(\cA)$ we have
\[\textstyle
\Re\langle\rho,A\rangle
=\Re\langle a(\rho,\dots,\rho),A\rangle
=\Re\langle a(\gamma(\rho),\dots,\gamma^k(\rho)),A\rangle
=\Re\langle\sigma,A\rangle\,,
\]
hence $\pi(\rho)=\sigma$. Regarding the inclusion ``$\subset$'',
let $\rho_i\in\cK$, $i=1,\dots,k$. Let 
$\sigma:=a(\gamma(\rho_1),\dots,\gamma^k(\rho_k))$ and 
$\tau:=a(\rho_1,\dots,\rho_k)$. Then for all $A\in\cH(\cA)$
\[\textstyle
\Re\langle\sigma,A\rangle
=\Re\langle a(\gamma(\rho_1),\dots,\gamma^k(\rho_k)),A\rangle
=\Re\langle a(\rho_1,\dots,\rho_k),A\rangle
=\Re\langle\tau,A\rangle
\]
holds. If $\sigma\in\cA$, then $\sigma=\pi(\tau)$ follows. 
As $\cK$ is convex, $\tau\in\cK$ holds and we obtain $\sigma\in\pi(\cK)$.
\end{proof}

\end{document}
